% Part: computability
% Chapter: computability-theory
% Section: fixed-point-thm

\documentclass[../../../include/open-logic-section]{subfiles}

\begin{document}

\olfileid{cmp}{thy}{fix}
\olsection{مبرهنة النقطة الثابتة}

لننظر مرة أخرى في مسألة التوقف. سنكتب مؤقتًا
$\gn{\cfind{x}(y)}$ بدلًا من $\tuple{x,y}$؛ ويمكن أن ننظر إلى هذا
بوصفه «اسمًا» لقيمة $\cfind{x}(y)$. وباستعمال هذا الترميز يمكننا
إعادة صياغة أحد براهيننا على أن مسألة التوقف غير قابلة للقرار.

السؤال: هل توجد دالة قابلة للحساب $h$ لها الخاصية الآتية؟ لكل $x$
و$y$،
\[
h(\gn{\cfind{x}(y)}) =
\begin{cases}
1 & \text{إذا كانت $\cfind{x}(y) \fdefined$} \\
0 & \text{خلاف ذلك.}
\end{cases}
\]

الجواب: لا؛ وإلا لكانت الدالة الجزئية
\[
g(x) \simeq
\begin{cases}
0 & \text{إذا كان $h(\gn{\cfind{x}(x)}) = 0$} \\
\text{غير معرّفة} & \text{خلاف ذلك}
\end{cases}
\]
قابلة للحساب، ومن ثم لكان لها فهرس ما~$e$. لكن عندئذ يكون
\[
\cfind{e}(e) \simeq
\begin{cases}
0 & \text{إذا كان $h(\gn{\cfind{e}(e)}) = 0$} \\
\text{غير معرّفة} & \text{خلاف ذلك،}
\end{cases}
\]
فتكون $\cfind{e}(e)$ معرّفة إذا وفقط إذا لم تكن معرّفة، وهذا تناقض.

تأمل الآن المعادلة التي تتضمن $\cfind{e}$. ففيها إحالة ذاتية بمعنى
ما: لقد رتبنا أن تعتمد قيمة $\cfind{e}(e)$، بطريقة معينة، على
$\gn{\cfind{e}(e)}$. وتقول مبرهنة النقطة الثابتة إنه \emph{يمكننا}
فعل ذلك بوجه عام، لا لمجرد إقامة التناقضات.

تعرض \olref{lem:fixed-equiv} صيغتين متكافئتين لمبرهنة النقطة
الثابتة. ومن الوجهة المنطقية ينتج تكافؤ العبارتين من كونهما صادقتين؛
لكن مقصودنا الحقيقي هو أن كلًا منهما تنتج من الأخرى مباشرة، بحيث
يمكن اتخاذهما صيغتين بديلتين للمبرهنة نفسها.

\begin{lem}
\ollabel{lem:fixed-equiv}
العبارتان الآتيتان متكافئتان:
\begin{enumerate}
\item لكل دالة جزئية قابلة للحساب $g(x,y)$ فهرس~$e$ بحيث، لكل~$y$،
\[
\cfind{e}(y) \simeq g(e,y).
\]
\item لكل دالة قابلة للحساب~$f(x)$ فهرس~$e$ بحيث، لكل~$y$،
\[
\cfind{e}(y) \simeq \cfind{f(e)}(y).
\]
\end{enumerate}
\end{lem}

\begin{proof}
$(1) \Rightarrow (2)$: إذا أعطيت $f$ فعرّف $g$ بالمعادلة
$g(x,y) \simeq \fn{Un}(f(x),y)$. واستعمل (1) لتحصل على فهرس~$e$
بحيث، لكل~$y$،
\begin{align*}
\cfind{e}(y) & = \fn{Un}(f(e),y) \\
& = \cfind{f(e)}(y).
\end{align*}

$(2) \Rightarrow (1)$: إذا أعطيت $g$ فاستعمل مبرهنة $s$-$m$-$n$
لتحصل على $f$ بحيث، لكل $x$ و~$y$،
$\cfind{f(x)}(y) \simeq g(x,y)$. ثم استعمل (2) لتحصل على فهرس~$e$
بحيث
\begin{align*}
\cfind{e}(y) & = \cfind{f(e)}(y) \\
& = g(e,y).
\end{align*}
وبذلك يتم البرهان.
\end{proof}

\begin{explain}
قبل أن نبين صدق العبارة (1)، ومن ثم صدق (2) أيضًا، تأمل مدى غرابتها.
انظر إلى $e$ على أنه برنامج حاسوب؛ تقول العبارة (1) إنك، إذا أعطيت
أي دالة جزئية قابلة للحساب $g(x,y)$، تستطيع إيجاد برنامج حاسوب $e$
يحسب $g_e(y) \simeq g(e,y)$. وبعبارة أخرى، تستطيع إيجاد برنامج يحسب
دالة تشير إلى البرنامج نفسه.
\end{explain}

\begin{thm}
العبارتان الواردتان في \olref{lem:fixed-equiv} صادقتان. وعلى وجه
التحديد، لكل دالة جزئية قابلة للحساب $g(x,y)$ فهرس~$e$ بحيث،
لكل~$y$،
\[
\cfind{e}(y) \simeq g(e,y).
\]
\end{thm}

\begin{proof}
مقومات البرهان كامنة بالفعل في المناقشة السابقة لمسألة التوقف. لتكن
$\fn{diag}(x)$ دالة قابلة للحساب تعيد، لكل $x$، فهرسًا للدالة
$f_x(y) \simeq \cfind{x}[2](x,y)$، أي
\[
\cfind{\fn{diag}(x)}(y) \simeq \cfind{x}[2](x,y).
\]
فكر في $\fn{diag}$ بوصفها دالة تحول برنامجًا لدالة ثنائية إلى برنامج
لدالة أحادية، وذلك بتثبيت البرنامج الأصلي وسيطًا أول لها. ويمكن تعريف
الدالة $\fn{diag}$ صوريًا كما يأتي: عرّف أولًا $s$ بالمعادلة
\[
s(x,y) \simeq \fn{Un}^2(x,x,y),
\]
حيث $\fn{Un}^2$ دالة ثلاثية شاملة للدوال الجزئية الثنائية القابلة
للحساب. ثم تتيح لنا مبرهنة $s$-$m$-$n$ إيجاد دالة عودية بدائية
$\fn{diag}$ تحقق
\[
\cfind{\fn{diag}(x)}(y) \simeq s(x,y).
\]

والآن عرّف الدالة $l$ بالمعادلة
\[
l(x,y) \simeq g(\fn{diag}(x),y).
\]
ولتكن $\gn{l}$ فهرسًا لـ$l$. وأخيرًا، ضع
$e=\fn{diag}(\gn{l})$. عندئذ، لكل $y$، لدينا
\begin{align*}
\cfind{e}(y) & \simeq \cfind{\fn{diag}(\gn{l})}(y) \\
& \simeq \cfind{\gn{l}}[2](\gn{l}, y) \\
& \simeq l(\gn{l}, y) \\
& \simeq g(\fn{diag}(\gn{l}),y) \\
& \simeq g(e, y),
\end{align*}
كما هو مطلوب.
\end{proof}

\begin{explain}
ما الذي يحدث هنا؟ افترض أنك كُلّفت بكتابة برنامج حاسوب يطبع نفسه.
وافترض أيضًا أنك تعمل بلغة برمجة لها مكتبة غنية وغريبة من دوال
السلاسل. ولنفترض، على وجه الخصوص، أن في لغتك دالة $\fn{diag}$ تعمل
كما يأتي: إذا أعطيت سلسلة دخل~$s$، حددت $\fn{diag}$ كل موضع يظهر فيه
الرمز `x' في~$s$ واستبدلت به نسخة مقتبسة من السلسلة الأصلية. فمثلًا،
إذا أعطيت السلسلة
\begin{quote}
\begin{verbatim}
hello x world
\end{verbatim}
\end{quote}
دخلًا، أعادت الدالة
\begin{quote}
\begin{verbatim}
hello 'hello x world' world
\end{verbatim}
\end{quote}
خرجًا. وعندئذ يسهل كتابة البرنامج المطلوب؛ ويمكنك التحقق من أن
\begin{quote}
\begin{verbatim}
print(diag('print(diag(x))'))
\end{verbatim}
\end{quote}
يفي بالغرض. وتنجح الفكرة نفسها في لغات أكثر شيوعًا مثل C++ وJava،
وإن كان تنفيذها أكثر تعقيدًا.

لم يبق بيننا وبين برهان مبرهنة النقطة الثابتة إلا خطوتان. افترض أن
صيغة معدلة من دالة الطباعة $\fn{print}(x,y)$ تقبل سلسلة $x$ ووسيطة
عددية أخرى $y$، وتطبع السلسلة $x$ عددًا من المرات مقداره $y$.
عندئذ يكون «البرنامج»
\begin{quote}
\begin{verbatim}
getinput(y);
print(diag('getinput(y); print(diag(x), y)'), y)
\end{verbatim}
\end{quote}
برنامجًا يطبع نفسه $y$ مرة عند الدخل $y$. وإذا استبدلنا بهيكل
$\fn{getinput}$---$\fn{print}$---$\fn{diag}$ دالة اعتباطية
$g(x,y)$ حصلنا على
\begin{quote}
\begin{verbatim}
g(diag('g(diag(x), y)'), y)
\end{verbatim}
\end{quote}
وهو برنامج يشغّل، عند الدخل $y$، الدالة $g$ على البرنامج نفسه وعلى
$y$. وإذا فهمنا «الاقتباس» على أنه «استعمال فهرس لـ» حصلنا على
البرهان السابق.

لا بأس الآن إن شئت أن تنظر إلى البرهان بوصفه حيلة صورية أو ضربًا من
السحر. لكن ينبغي أن تستطيع إعادة بناء تفاصيل الحجة السابقة. وحين
نبرهن مبرهنات عدم الاكتمال (و«مبرهنة النقطة الثابتة» المتصلة بها)
سنناقش طرائق أخرى لفهم سبب نجاحها.
\end{explain}

\begin{tagblock}{lambda}
\begin{digress}
يمكن استعمال الفكرة نفسها للحصول على مؤلّف «نقطة ثابتة». افترض أن
لديك حد لامبدا $g$، وأنك تريد حدًا آخر $k$ له الخاصية أن $k$ مكافئ
$\beta$ للحد $gk$. عرّف الحدين
\[
\fn{diag}(x) = xx
\]
و
\[
l(x) = g(\fn{diag}(x))
\]
وفق اصطلاحاتنا الترميزية؛ أي إن $l$ هو الحد
$\lambd[x][g(xx)]$. ولْيكن $k$ الحد $ll$. عندئذ لدينا
\begin{align*}
k & = (\lambd[x][g(xx)])(\lambd[x][g(xx)]) \\
& \red  g((\lambd[x][g(xx)])(\lambd[x][g(xx)])) \\
& = gk.
\end{align*}
فإذا أخذنا
\[
Y = \lambd[g][((\lambd[x][g(xx)])(\lambd[x][g(xx)]))]
\]
اختزل $Yg$ و$g(Yg)$ إلى حد مشترك؛ ومن ثم
$Yg \equiv_\beta g(Yg)$. ويعرف هذا باسم «مؤلّف كاري». أما إذا أخذنا
بدلًا منه
\[
Y = (\lambd[xg][g(xxg)])(\lambd[xg][g(xxg)])
\]
فإن $Yg$ يختزل في الواقع إلى $g(Yg)$، وهي عبارة أقوى. وتعرف هذه
الصيغة الأخيرة من $Y$ باسم «مؤلّف تورنغ».
\end{digress}
\end{tagblock}

\end{document}
